\documentclass[11pt]{article}
\usepackage[a4paper,margin=2.5cm]{geometry}
\usepackage{amsmath,amssymb}
\usepackage{booktabs}
\usepackage{longtable}
\usepackage[hidelinks]{hyperref}

\title{Diets Used in Kaun et al. (2007):\\DEB-Consistent Formulation}
\author{}
\date{}

\begin{document}
\maketitle

This note extracts and formalizes the larval diets used in Kaun et al. (2007) and translates them into DEB-consistent substrate descriptions, ending with a single derived quantity---the environmental energy density of food $\varepsilon_X$ (J\,cm$^{-3}$)---that can be used to parameterize functional response and assimilation in DEB models.

Throughout, we explicitly separate
\begin{itemize}
  \item the DEB \emph{food density} $X$ (substrate density, e.g. C-mol\,cm$^{-3}$),
  \item the \emph{chemical potential} $\mu_X$ (J\,C-mol$^{-1}$), and
  \item the derived \emph{environmental energy density}
  \begin{equation}
    \varepsilon_X \equiv \mu_X X, \qquad [\varepsilon_X] = \mathrm{J\,cm^{-3}}.
  \end{equation}
\end{itemize}

This quantity is not a new DEB state variable; it is a defined product that is useful when food recipes are given by mass or molarity.

\section{Diets detected in Kaun et al. (2007)}

Kaun et al. use three classes of larval food substrates, depending on the experiment:
\begin{enumerate}
  \item standard rearing medium (yeast--sucrose--agar), sometimes diluted (100\%, 75\%, 50\%, 25\%, 15\%),
  \item yeast paste (used in feeding and absorption assays), and
  \item defined sugar--agarose substrates (glucose--agarose and fructose--agarose).
\end{enumerate}

Only the standard medium and yeast paste are nutritionally complete. The sugar--agarose substrates are single-compound carbon sources, used to probe feeding behaviour rather than growth.

\section{Notational conventions (important)}

To avoid confusion, we enforce the following conventions:
\begin{itemize}
  \item $X$: DEB food density (substrate), units chosen consistently with $K$ (e.g. C-mol\,cm$^{-3}$),
  \item $\mu_X$: chemical potential of food (J\,C-mol$^{-1}$),
  \item $\varepsilon_X = \mu_X X$: energy per volume of substrate (J\,cm$^{-3}$).
\end{itemize}

When the paper specifies recipes in g\,L$^{-1}$ or mol\,L$^{-1}$, we first compute $\varepsilon_X$ via mass concentrations, using
\begin{equation}
  \varepsilon_X = \sum_i \rho_i e_i,
\end{equation}
where $\rho_i$ is the mass concentration of ingredient $i$ (g\,cm$^{-3}$) and $e_i$ is the gross energy content of ingredient $i$ (J\,g$^{-1}$). Conversion to $X$ (e.g. C-mol\,cm$^{-3}$) can be done later if needed.

\section{Diet 1: Standard rearing medium}

\subsection{Composition (per litre of water)}

A typical standard medium recipe is summarised in Table~\ref{tab:standard-composition}.

\begin{table}[h]
\centering
\begin{tabular}{lll}
\toprule
Component & Amount (g\,L$^{-1}$) & Role \\
\midrule
Baker's yeast & 50 & protein, lipid, carbohydrate \\
Sucrose & 100 & carbohydrate \\
Agar & 16 & gelling agent (non-assimilable) \\
KPO$_4$ & 0.1 & mineral \\
NaK-tartrate & 8 & mineral \\
NaCl & 0.5 & mineral \\
MgCl$_2$ & 0.5 & mineral \\
Fe$_2$(SO$_4$)$_3$ & 0.5 & mineral \\
Water & balance & solvent \\
\bottomrule
\end{tabular}
\caption{Composition of the standard larval medium per litre of water in Kaun et al. (2007). Only yeast and sucrose contribute significantly to assimilable energy.}
\label{tab:standard-composition}
\end{table}

Food quality manipulations (25\%, 15\%, etc.) scale yeast and sucrose together.

\subsection{Energy density of the standard medium}

Assume an aqueous gel with bulk density $\rho_{\text{medium}} \approx 1\,\mathrm{g\,cm^{-3}}$ and negligible assimilable energy from agar. Take typical energy per mass values
\begin{align}
  e_Y &\approx 21\times10^{3}\,\mathrm{J\,g^{-1}} && \text{(yeast)},\\
  e_S &\approx 16.5\times10^{3}\,\mathrm{J\,g^{-1}} && \text{(sucrose)}.
\end{align}

Energy per litre from yeast and sucrose is then
\begin{align}
  E_{\text{yeast}} &= 50\times 21 = 1050\,\mathrm{kJ},\\
  E_{\text{sucrose}} &= 100\times 16.5 = 1650\,\mathrm{kJ},\\
  E_{\text{total}} &= 2700\,\mathrm{kJ\,L^{-1}}.
\end{align}

Thus the energy density of the standard medium is
\begin{equation}
  \varepsilon_X^{\text{standard}} = \frac{2.7\times10^{6}}{1000} \approx 2.7\times10^{3}\,\mathrm{J\,cm^{-3}}.
\end{equation}

For diluted food with quality factor $q$ (e.g. 0.75, 0.50, 0.25, 0.15),
\begin{equation}
  \varepsilon_X(q) = q\,\varepsilon_X^{\text{standard}}.
\end{equation}

\section{Diet 2: Yeast paste}

\subsection{Composition}

The yeast paste is prepared as 2:1 water:yeast (w/w). The mass fraction of yeast is
\begin{equation}
  w_Y = \frac{1}{3}.
\end{equation}
Assuming slurry density $\approx 1\,\mathrm{g\,cm^{-3}}$, the yeast mass concentration is
\begin{equation}
  \rho_Y = w_Y\,\rho_{\text{paste}} = 0.333\,\mathrm{g\,cm^{-3}}.
\end{equation}

\subsection{Energy density}

With yeast energy per mass $e_Y = 21\times10^{3}\,\mathrm{J\,g^{-1}}$ we obtain
\begin{equation}
  \varepsilon_{X,\text{yeast paste}} = \rho_Y e_Y = 0.333\times 21\times10^{3} \approx 7.0\times10^{3}\,\mathrm{J\,cm^{-3}}.
\end{equation}

Unit check: $\mathrm{g\,cm^{-3}}\times \mathrm{J\,g^{-1}} = \mathrm{J\,cm^{-3}}$.

In the $^{14}$C-glucose assays of Kaun et al. (2007), this yeast-paste energy density is the dominant energetic term. The labelled glucose is treated as a tracer signal for pathway-specific ingestion/retention; when no tracer mass/volume density is reported, the tracer is not introduced as an additional independent dietary energy input in DEB energetics.

\section{Diet 3: Sugar--agarose substrates}

These substrates are not nutritionally complete and should be treated as single-compound DEB foods.

\subsection{Glucose--agarose}

Recipe: 10\% glucose (w/v), 2.3\% agarose. The glucose mass concentration is
\begin{equation}
  \rho_G = 0.10\,\mathrm{g\,cm^{-3}}.
\end{equation}
With glucose energy per mass $e_G = 15.6\times10^{3}\,\mathrm{J\,g^{-1}}$,
\begin{equation}
  \varepsilon_{X,\text{glucose}} = \rho_G e_G = 0.10\times 15.6\times10^{3} = 1.56\times10^{3}\,\mathrm{J\,cm^{-3}}.
\end{equation}

\subsection{Fructose--agarose}

Recipe: 2 mol\,L$^{-1}$ fructose, 1\% agarose. With fructose molar mass $M_F = 180\,\mathrm{g\,mol^{-1}}$, the mass concentration is
\begin{equation}
  \rho_F = \frac{2\times 180}{1000} = 0.36\,\mathrm{g\,cm^{-3}}.
\end{equation}
With carbohydrate energy per mass $e_F \approx 15.6\times10^{3}\,\mathrm{J\,g^{-1}}$,
\begin{equation}
  \varepsilon_{X,\text{fructose}} = 0.36\times 15.6\times10^{3} \approx 5.6\times10^{3}\,\mathrm{J\,cm^{-3}}.
\end{equation}

\section{Summary: single food energy densities}

\begin{table}[h]
\centering
\begin{tabular}{lll}
\toprule
Diet & $\varepsilon_X$ (J\,cm$^{-3}$) & Notes \\
\midrule
Standard medium (100\%) & $2.7\times10^{3}$ & complete diet \\
Standard medium (quality $q$) & $q\times 2.7\times10^{3}$ & dilution series \\
Yeast paste & $7.0\times10^{3}$ & high-quality reference \\
Glucose--agarose & $1.56\times10^{3}$ & carbon-only test \\
Fructose--agarose & $5.6\times10^{3}$ & carbon-only test \\
\bottomrule
\end{tabular}
\caption{Summary of energy densities for Kaun et al. larval diets.}
\label{tab:summary}
\end{table}

\end{document}
